\setcounter{secnumdepth}{-1}

\chapter{Introducción}\footnote{El formato y la estructura documental de esta memoria han sido desarrollados utilizando la plantilla oficial de proyectos de la ULB~\cite{ulb_template}.}

En el contexto de la \textbf{transformación digital de las relaciones interpersonales}, las plataformas de citas y aplicaciones de \textit{matching} han adquirido un papel central en la forma en que las personas establecen conexiones afectivas. En la actualidad, estas plataformas gestionan enormes volúmenes de información sobre preferencias, rasgos de personalidad, características demográficas y patrones de comportamiento. El objetivo fundamental de este ecosistema digital es procesar dicha información para facilitar la formación de \textbf{relaciones compatibles y potencialmente duraderas}, actuando como intermediarios algorítmicos en el tejido social.\\

Tradicionalmente, los sistemas de recomendación en este ámbito se han basado en filtros simples, reglas heurísticas o mecanismos poco transparentes que priorizan coincidencias superficiales (como la edad, la localización geográfica o los intereses básicos). Sin embargo, la compatibilidad relacional es un \textbf{fenómeno inherentemente complejo y multidimensional}. El éxito de una relación no se define por la simple superposición de aficiones, sino que depende de una intrincada red de factores psicológicos, valores personales, estilo de vida, ambiciones profesionales y patrones emocionales. La incapacidad de los sistemas tradicionales para capturar esta profundidad suele derivar en recomendaciones irrelevantes, generando frustración y fatiga en el usuario.\\

En este contexto, surge la necesidad crítica de diseñar \textbf{sistemas de inteligencia artificial} capaces de modelar de forma cuantitativa y estructurada la compatibilidad entre individuos. Al integrar múltiples dimensiones personales, estos sistemas pueden descubrir patrones latentes y generar predicciones que sean \textbf{consistentes, reproducibles y explicables}. En el delicado ámbito de las relaciones humanas, un sistema de este tipo no solo debe alcanzar un alto rendimiento técnico o estadístico, sino también garantizar su interpretabilidad. Resulta imperativo mantener la coherencia con los principios del dominio relacional, evitando a toda costa funcionar como una "caja negra" que tome decisiones opacas o arbitrarias.\\

El presente proyecto se fundamenta en el dataset \textit{“Cupid’s Algorithm”} ~\cite{dataset_cupid}, un conjunto de datos que contiene información estructurada y enriquecida sobre pares de individuos. Este \textit{dataset} abarca desde variables demográficas y profesionales hasta métricas psicométricas profundas, como los rasgos de personalidad basados en el modelo \textbf{\textit{Big Five}}~\cite{rakic2024bigfive} y los lenguajes del amor, acompañados de indicadores históricos de compatibilidad y duración potencial de la relación. La riqueza de esta información permite abordar el reto desde una \textbf{perspectiva supervisada}, formalizando la afinidad relacional como una variable matemática susceptible de ser modelada mediante técnicas avanzadas de aprendizaje automático.\\

A partir de este marco conceptual, el proyecto no se limita al mero entrenamiento de un modelo predictivo aislado en un entorno académico, sino que adopta una visión integral del \textbf{sistema de inteligencia artificial como producto}. Esto implica recorrer, de manera rigurosa, todas las fases del ciclo de vida de un sistema de \textit{Machine Learning} (MLOps): desde la definición formal del problema y la preparación de los datos, pasando por el modelado y la evaluación, hasta culminar en su \textbf{integración como servicio de software} y su correspondiente análisis de viabilidad técnica y económica.\\

De este modo, el objetivo principal del proyecto es diseñar e implementar un \textbf{sistema inteligente} capaz de predecir la compatibilidad entre dos individuos de forma completamente objetiva y explicable. Un sistema concebido desde su inicio para ser un componente integrable en una plataforma digital de relaciones, y evaluado no solo a través de métricas de precisión técnica, sino también por su \textbf{potencial impacto organizativo, económico y social}.